% =============================================================
% Chapter 9: Conclusion & Future Work
% =============================================================

\newpage
\vspace*{\fill}
\begin{center}
\textbf{\Huge{Chapter 9}}\\
\vspace{5mm}
\textbf{\Huge{Conclusion \& Future Work}}
\end{center}
$1
\begin{center}
\thepage
\end{center}
$2

\titleformat{\chapter}[block]{\normalfont}{}{0pt}{}
\titlespacing{\chapter}{0pt}{-2em}{0pt}
\chapter[Conclusion and Future Work]{}
\titleformat{\chapter}[display]{\vspace{1cm}\normalfont\Large\bfseries\centering}{\chaptertitlename\ \thechapter}{20pt}{\vspace{0pt}}
\titlespacing*{\chapter}{0pt}{-50pt}{20pt}

\section{Conclusion}
The creation of \textbf{AyurSpace} is a key step in the digital preservation and democratization of Ayurvedic knowledge. By effectively executing a hybrid neuro-symbolic architecture, this framework deals with the major problems facing current plant identification tools---specifically, the absence of medicinal context and the risk of generative hallucinations. Our experimental findings, producing a 96.4\% accuracy in taxonomy and a 99.1\% compliance in contextual generative reasoning, confirm the effectiveness of decoupling visual recognition from semantic analysis. This dependability is extremely important in the health field, where misidentification can have toxicological effects. Additionally, the algorithmic formalization of the \textit{Dosha} assessment transforms subjective clinical intuition into a replicable digital logic, enabling users to make smart health decisions based on ancient wisdom merged with robust modern mobile processing.

\section{Future Scope}
Future research and scale expansions will focus on three main development branches:
\begin{enumerate}
    \item \textbf{Edge AI Optimization}: Utilizing quantized MobileNet models (int8 quantization) on the mobile device natively. This will allow for the complete offline visual identification of plants, resolving accessibility issues for users in remote rural locations suffering from weak network connectivity.
    \item \textbf{Augmented Reality (AR) Interfaces}: Harnessing ARCore frameworks to overlay medicinal properties, alerts, and detailed plant structures directly onto the camera viewfinder. This transition to augmented realities will cultivate deeply immersive, real-time educational experiences for students and the general public.
    \item \textbf{Telemedicine Integration}: Launching a "Vaidya Connect" application component. By piping users' generated Dosha baselines and scan histories into professional Ayurvedic dashboards, we can create an end-to-end telemedicine ecosystem that blends self-evaluation algorithms with direct oversight from licensed Ayurvedic doctors.
\end{enumerate}
